\subsection{数据工程：所谓“干净互联网”是加工结果}

slides 12--15 是本讲最容易被一句“用 Common Crawl 训练”掩盖的部分。真实流水线从 WARC 抓取档案开始，要经历 HTML 正文抽取、NSFW/PII 等风险过滤、URL/文档/行级去重、启发式质量过滤、模型质量评分和 domain mixture 重加权。每一步都在改变训练分布。midtraining 则在预训练后用更小但更有针对性的语料继续训练，用于提高科学、代码、多语言、长上下文、格式遵循或 reasoning 等特性。

\lecturefigure{slides-images/slide-012.png}{官方 slide 12：从 Common Crawl 到数据 mixture 的七步预训练数据流水线}
\lecturefigure{slides-images/slide-013.png}{官方 slide 13：FineWeb 各级过滤如何逐步缩小网页集合}
\lecturefigure{slides-images/slide-014.png}{官方 slide 14：midtraining 的领域、长度、格式与质量适配}
\lecturefigure{slides-images/slide-015.png}{官方 slide 15：公开语料规模、研究问题与数据保密动机}

\begin{knowledgebox}{读图：过滤漏斗同时改变规模与分布}
slide 13 右侧的 200B、36B、20B、15B 不是“越少越好”的排行榜，而是网页集合经过语言、重复、长度、模板和质量规则后的数量变化。读者应追问每层删掉了什么、对哪些语言或网站更不利、是否把稀有但有价值的内容一起删除。slide 15 列出的 C4、The Pile、Dolma、FineWeb 也不能只按 token 数比较；采集时间、许可、语言比例、去重粒度和分类体系都会影响模型学到的行为。
\end{knowledgebox}

\begin{warningbox}{数据质量不是单一“清洁分数”}
过滤器通常同时承担安全、法律、语言和质量目标，这些目标可能冲突。强力去重可降低 memorization，却可能删除多来源共同确认的事实；Wikipedia-like classifier 可提升百科风格，却可能压低论坛、口语和低资源语言；高质量 annealing 可改善终局指标，却不证明原始大规模语料可以被省略。\textbf{课堂提示：}讲者把数据称为 practical LLM 的关键，意在强调流水线和 mixture 的决定性，而不是暗示存在一个放之四海皆准的“完美语料”。
\end{warningbox}

从 Agent 角度看，这一段还有更深的连接：预训练分布决定模型最初会哪些工具、熟悉哪些接口、默认采用什么写作和推理习惯；midtraining 决定这些先验能否适应长上下文、代码或多语言环境。后训练可以校正行为，却很难低成本补回基础模型从未获得的广泛知识。因此数据治理既是模型能力问题，也是后续 Agent 环境覆盖率问题。

